Every figure quoted below is generated from the run logs by
\texttt{src/analyse.py}; none is typed by hand.

\subsection{Flow queueing dominates}

Table~\ref{tab:steady} gives all ten conditions. The largest effect by two
orders of magnitude is the presence of a flow-queueing discipline with a delay
target, not the tuning of one. Under the steady workload an unmanaged FIFO
reaches \PfifoProbeRttPninetyfiveSteady\,ms p95 RTT against
\FqCodelProbeRttPninetyfiveSteady\,ms for \texttt{fq\_codel} --- a 99.0\%
reduction ($p<0.001$) --- at goodput that is statistically indistinguishable
(\PfifoGoodputSteady\ vs \FqCodelGoodputSteady\,Mbps, $p=0.20$). The staged
workload gives the same picture, and there \texttt{pfifo} additionally
\emph{loses} 10.2\% goodput and 1.4\% fairness ($p=0.015$, $p=0.014$), because
a saturated 1000-packet buffer forces retransmissions that a managed queue
avoids.

Fair queueing without a delay target is not sufficient. SFQ isolates flows but
applies no AQM, and its bulk queues sit at its 127-packet limit throughout:
its sparse probe reads \SfqProbeRttMeanSteady\,ms while the bulk flows sharing
the link read \SfqBulkRttMeanSteady\,ms. The AQM and the scheduler do distinct
jobs and both are necessary.

\subsection{Sparse and bulk latency are not interchangeable}

Figure~\ref{fig:fig02_latency_sparse_vs_bulk_steady} plots the two
measurements side by side. The ratio between them ranges from 1.1$\times$
(CoDel, a single queue) to 9.6$\times$ (SFQ), because \texttt{fq\_codel}'s
new-flow heuristic gives a sparse probe preferential service by design.

A single-probe measurement can therefore invert a ranking. By
sparse probe alone SFQ (\SfqProbeRttMeanSteady\,ms) appears to beat CoDel
(\CodelProbeRttMeanSteady\,ms); by the latency the bulk flows actually
experience, SFQ is \SfqBulkRttMeanSteady\,ms against CoDel's
\CodelBulkRttMeanSteady\,ms. We report both throughout, and we suggest any AQM
evaluation should.

\subsection{The controller's computational cost is negligible}

The sham condition runs the controller at its full polling cadence and applies
no parameter change. Against static \texttt{fq\_codel} it is indistinguishable
on every metric in both workloads, p95 probe RTT $p=0.29$ (steady) and
$p=0.24$ (staged), backlog $p=0.94$ and $p=0.76$.

This matters for interpreting the next subsection. Without it, any difference
between static and adapted \texttt{fq\_codel} would confound the controller's
CPU cost with its control decisions, and we could not attribute an effect to
either. Having established the cost is nil, differences that remain are
attributable to the decisions.

\subsection{Adaptation: a small effect under steady load, none under change}

Under the \textbf{steady} workload the adaptation produces a real but small
improvement: mean backlog falls \FqCodelBacklogSteady\ from
\FqCodelAcapeBacklogSteady\ packets, a 9.6\% reduction ($p=0.050$), and p95
probe RTT falls 0.5\% ($p=0.047$). Goodput and fairness are unchanged. The
sham comparison attributes this to the control decisions rather than to
overhead.

Under the \textbf{staged} workload, where the flow count changes twice during
the run, the adaptation produces \emph{no} improvement on any metric. Mean
probe RTT is 1.6\% \emph{worse} ($p=0.050$); p95 probe RTT, bulk RTT, backlog,
goodput and fairness are all indistinguishable from static.

That direction is worth stating plainly. The workload designed to
give the controller something to respond to is the one in which it helps least.
Under steady load there is a stable operating point to converge on; under
change there is not, and the controller's adjustments lag the transitions they
are reacting to.

\subsection{An off-the-shelf alternative outperforms the adapted system}

CAKE, configured with no controller and a single \texttt{besteffort} argument,
achieves \CakeProbeRttPninetyfiveSteady\,ms p95 probe RTT under the steady
workload against \FqCodelAcapeProbeRttPninetyfiveSteady\,ms for adapted
\texttt{fq\_codel} ($p<0.001$), and \CakeBulkRttMeanSteady\,ms bulk RTT against
\FqCodelAcapeBulkRttMeanSteady\,ms.

The margin by which CAKE beats the adapted system is larger than the margin by
which the adaptation beats the static defaults. For a practitioner seeking
lower latency on this class of link, changing queue discipline is a larger and
far simpler win than tuning \texttt{fq\_codel}'s parameters.

\subsection{Predictive control never engaged}

No adjustment in any run was tagged predictive. The cause is structural rather
than a coding fault, and it survives the correction described in
Section~\ref{sec:pitfalls}: an adjustment fires only once the regime has been
unchanged for five consecutive samples, and the trajectory estimate is derived
from the gradient over that same window, which by then is below the threshold
that would classify it as WORSENING or RECOVERING. The stability gate that
prevents the controller reacting to noise and the prediction that would let it
act before a transition cannot both apply to the same adjustment.

We therefore withdraw the predictive-control contribution rather than claim it.
Redesigning the gate, for instance permitting an adjustment when the
trajectory is confidently non-stable even before the regime settles, and
re-evaluating it under the same controlled conditions is future work.

\subsection{Where adaptation does help: short paths}

CoDel's \texttt{target} is a delay budget expressed in absolute units, while
its authors argue the default is appropriate relative to path round-trip time.
Sweeping the base RTT over 5, 20, 80 and 200\,ms tests that directly.

Static \texttt{fq\_codel} holds median queueing delay between 1.96 and
5.00\,ms across the entire 40-fold span, bounded by its 5\,ms target. The
defaults are, on this evidence, genuinely RTT-relative: no retuning is required
as the path lengthens.

Adaptation is nonetheless not uniformly inert. Measured on the RTT the bulk
flows themselves experience:

\begin{center}
\begin{tabular}{rrrr}
\toprule
\textbf{Base RTT} & \textbf{Static} & \textbf{Adapted} & \textbf{Change} \\
\midrule
5\,ms   & $27.68 \pm 0.14$  & $22.58 \pm 0.15$  & $-18.4\%$\textsuperscript{*} \\
20\,ms  & $36.39 \pm 0.08$  & $35.87 \pm 0.80$  & $-1.4\%$ \\
80\,ms  & $89.30 \pm 0.37$  & $89.86 \pm 0.70$  & $+0.6\%$ \\
200\,ms & $208.00 \pm 0.76$ & $207.47 \pm 0.74$ & $-0.3\%$ \\
\bottomrule
\end{tabular}
\end{center}

\noindent\textsuperscript{*}$p < 0.0001$; the remaining three are not
distinguishable ($p = 0.38$, $0.31$, $0.43$).

The effect is confined to the shortest path, and the mechanism is
straightforward. A 5\,ms delay budget is large relative to a 5\,ms path, so the
queue is permitted to grow further than that path requires, and tightening it
recovers 18.4\% of the bulk flows' round-trip time. Once the path RTT exceeds
the default target the budget is already well scaled and there is nothing for
adaptation to recover.

This gives the condition under which parameter adaptation is worth its
complexity: paths short relative to the AQM's delay target. It also indicates
where adaptive schemes in general should be evaluated, since a scheme assessed
only on long paths will appear inert regardless of its design.

\subsection{Sparse flows under a mixed workload}

Bulk-only traffic cannot exercise \texttt{fq\_codel}'s new-flow heuristic or
its \texttt{quantum} parameter, both of which exist to serve sparse
latency-sensitive flows. Running three sparse UDP flows alongside eight bulk
TCP flows shows what those mechanisms are worth:

\begin{center}
\begin{tabular}{lrr}
\toprule
\textbf{System} & \textbf{Sparse loss (\%)} & \textbf{Sparse jitter (ms)} \\
\midrule
\texttt{pfifo}, no AQM & 7.97 & 2.19 \\
FQ-PIE                 & 0.000 & 1.91 \\
CAKE                   & 0.007 & \textbf{1.36} \\
\texttt{fq\_codel}     & 0.002 & 2.00 \\
\bottomrule
\end{tabular}
\end{center}

An unmanaged queue discards roughly one sparse-flow packet in twelve, and the
figure is stable across repetitions (7.47\%, 7.98\%, 8.31\%). Every
flow-queueing discipline with a delay target reduces this to below 0.03\%. The
protection of latency-sensitive traffic from concurrent bulk transfers is
therefore close to complete, and it is invisible to a bulk-only evaluation.

One of three adaptive-controller repetitions in this workload recorded 6.31\%
sparse loss against 0.000\% and 0.029\% in the others. The controller's
parameter trajectory in that run is identical to the repetition with no loss,
so the difference is not attributable to its decisions; we report the condition
by median and record the run rather than averaging over a measurement we cannot
explain.
